\documentclass{article}
\usepackage{graphicx} % Required for inserting images
\usepackage{amsmath}
\usepackage{algorithmic}
\usepackage{algorithm}
\usepackage[polish]{babel}
\usepackage[utf8]{inputenc}
\usepackage[T1]{fontenc}
\usepackage{float}
\usepackage{listings}
\usepackage[table]{xcolor}
\usepackage{diagbox}
\usepackage{geometry}

\title{Algorytmy Równoległe - Laboratorium 5}
\author{Adam Staniszewski}

\lstset{
    language=bash,
    basicstyle=\ttfamily,
    keywordstyle=\color{blue},
    commentstyle=\color{gray},
    stringstyle=\color{orange},
    backgroundcolor=\color{lightgray!20},
    frame=single,
    captionpos=b
}

\begin{document}
\maketitle
\tableofcontents

\section{Instalacja narzędzia}
Do analizy użyjemy narzędzia \textbf{VTune}, które załadujemy na Aresie przy pomocy komendy:

\begin{lstlisting}
module load vtune
\end{lstlisting}

\section{Testy narzędzia}
W ramach testów uruchomimy algorytm na czterech rdzeniach przy użyciu narzędzia \textbf{VTune}. Algorytm będzie rozwiązywał problem o rozmiarze \textbf{16 na 16}.
Skrypt uruchamiamy następująco:
\begin{lstlisting}
vtune -collect hotspots -duration 30 -result-dir vtune_results 
python lab3.py
\end{lstlisting}
\vfill
Po wywołaniu komedy, VTune wyświetli w terminalu:
\begin{figure}[H]
    \centering
    \includegraphics[width=0.8\linewidth]{vtune_result1.png}
    \caption{VTune w terminalu}
    \label{fig:enter-label}
\end{figure}
Możemy zauważyć bardzo niską \textbf{efektywność wykorzystania CPU}. Prawdopodobnie jest to związane z:
\begin{itemize}
    \item niewielkim rozmiarem problemu,
    \item oczekiwaniami VTune'a odnośnie ilości CPU - widzimy, że oprogramowanie oczekuje użycia 48 procesorów,
\end{itemize}
Zasadne wydaje się zatem znaczne zwiększenie rozmiaru problemu i ilości rdzeni w ramach dalszych eksperymentów.
\noindent
W tym laboratorium mużyjemy \textbf{metody Gaussa-Seidla} do rozwiązania zadania.

\section{Analiza wydajności algorytmu dla różnych rozmiarów problemów}
Rozszerzmy wyżej wspomniane eksperymenty o inne wielkości problemów i dodatkowe rdzenie. Poprawmy także zawarty w nich błąd - rozmiary problemów powinny być kolejnymi \textbf{potęgami dwójki} większymi od 2.

\subsection{Nowy algorytm}
\begin{lstlisting}
import time
import numpy as np
import matplotlib.pyplot as plt
from mpi4py import MPI

def main():
    Lx: float = 10.0
    Ly: float = 10.0
    Nx: int = 64
    Ny: int = 64
    g: float = 1.0
    lambda_: float = 1.0
    omega: float = 1.8 

    dx = Lx / (Nx - 1)
    dy = Ly / (Ny - 1)

    comm = MPI.COMM_WORLD
    rank = comm.Get_rank()
    size = comm.Get_size()

    rows_per_process = Ny // size
    extra_rows = Ny % size
    local_rows = rows_per_process + 1 if rank < extra_rows else rows_per_process

    T_local = np.zeros((local_rows + 2, Nx)) 
    T_local[0, :] = 0 
    T_local[-1, :] = 0  

    tolerance: float = 1e-6
    error: float = 1.0

    while error > tolerance:
        error_local = 0.0

        for j in range(1, local_rows + 1):
            for i in range(1, Nx - 1):
                T_new = (T_local[j + 1, i] + T_local[j - 1, i] +
                         T_local[j, i + 1] + T_local[j, i - 1] - 
                         (g / lambda_) * (dx ** 2 + dy ** 2)) / 4.0
                T_local[j, i] = (1 - omega) * T_local[j, i] + omega * T_new
                error_local = max(error_local, abs(T_local[j, i] - T_new))

        reqs = []
        if rank > 0:
            reqs.append(comm.Irecv(T_local[0, :], source=rank - 1, tag=0))
            reqs.append(comm.Isend(T_local[1, :], dest=rank - 1, tag=1))
        if rank < size - 1:
            reqs.append(comm.Irecv(T_local[-1, :], source=rank + 1, tag=1))
            reqs.append(comm.Isend(T_local[-2, :], dest=rank + 1, tag=0))

        MPI.Request.Waitall(reqs)

        error = comm.allreduce(error_local, op=MPI.MAX)

    T_local_flat = T_local[1:-1, :].flatten() 
    counts = np.array(comm.gather(len(T_local_flat), root=0))

    if rank == 0:
        T_global_flat = np.empty(sum(counts), dtype=T_local.dtype)
    else:
        T_global_flat = None

    comm.Gatherv(T_local_flat, (T_global_flat, counts), root=0)

    if rank == 0:
        T_global = np.zeros((Ny, Nx))
        row_offset = 0
        for i in range(size):
            local_rows_i = rows_per_process + 1 if i < extra_rows 
            else rows_per_process
            T_global[row_offset:row_offset + local_rows_i, :] = \
                T_global_flat[row_offset * Nx:(row_offset
                + local_rows_i) * Nx].reshape(local_rows_i, Nx)
            row_offset += local_rows_i
        return T_global
    return None

if __name__ == "__main__":
    start = time.time()
    T_global = main()
    end = time.time()
    print(f"Execution Time: {(end - start):.7f} seconds")
\end{lstlisting}

\subsection{Rozmiar 32 na 32}
\begin{table}[ht]
\centering
\begin{tabular}{|l|r|r|r|}
\hline
\rowcolor{orange!40} 
Liczba rdzeni & Średni czas wykonania algorytmu (s) \\\hline
1 (implementacja sekwencyjna) & 2.91 \\\hline
2 & 1.49 \\\hline
4 & 0.28 \\\hline
8 & 0.26 \\\hline
16 & 0.22 \\\hline
32 & 0.23 \\\hline
48 & 0.24 \\\hline

\end{tabular}
\caption{\label{tab:widgets}Średni czas obliczeń dla problemu o rozmiarze 32 na 32.}
\end{table}

\begin{figure}[H]
    \centering
    \includegraphics[width=1\linewidth]{32_32_4.png}
    \caption{Rozmiar 32 na 32, 4 procesory.}
    \label{fig:enter-label}
\end{figure}

\begin{figure}[H]
    \centering
    \includegraphics[width=1\linewidth]{32_32_8.png}
    \caption{Rozmiar 32 na 32, 8 procesorów.}
    \label{Rozmiar 32 na 32, 8 procesorów.}
\end{figure}

\begin{figure}[H]
    \centering
    \includegraphics[width=1\linewidth]{32_32_16.png}
    \caption{Rozmiar 32 na 32, 16 procesorów.}
    \label{fig:enter-label}
\end{figure}

\begin{figure}[H]
    \centering
    \includegraphics[width=1\linewidth]{32_32_32.png}
    \caption{Rozmiar 32 na 32, 32 procesory.}
    \label{fig:enter-label}
\end{figure}

\begin{figure}[H]
    \centering
    \includegraphics[width=1\linewidth]{32_32_48.png}
    \caption{Rozmiar 32 na 32, 48 procesorów.}
    \label{fig:enter-label}
\end{figure}

\subsection{Rozmiar 64 na 64}
\begin{table}[ht]
\centering
\begin{tabular}{|l|r|r|r|}
\hline
\rowcolor{orange!40} 
Liczba rdzeni & Średni czas wykonania algorytmu (s) \\\hline
1 (implementacja sekwencyjna) & 38.14 \\\hline
2 & 22.01 \\\hline
4 & 4.97 \\\hline
8 & 4.56 \\\hline
16 & 4.31 \\\hline
32 & 4.78 \\\hline
48 & 45.16 \\\hline

\end{tabular}
\caption{\label{tab:widgets}Średni czas obliczeń dla problemu o rozmiarze 64 na 64.}
\end{table}

\begin{figure}[H]
    \centering
    \includegraphics[width=1\linewidth]{63_64_4.png}
    \caption{Rozmiar 64 na 64, 4 procesory.}
    \label{fig:enter-label}
\end{figure}

\begin{figure}[H]
    \centering
    \includegraphics[width=1\linewidth]{64_64_8.png}
    \caption{Rozmiar 64 na 64, 8 procesorów.}
    \label{fig:enter-label}
\end{figure}

\begin{figure}[H]
    \centering
    \includegraphics[width=1\linewidth]{64_64_16.png}
    \caption{Rozmiar 64 na 64, 16 procesorów.}
    \label{fig:enter-label}
\end{figure}

\begin{figure}[H]
    \centering
    \includegraphics[width=1\linewidth]{64_64_32.png}
    \caption{Rozmiar 64 na 64, 32 procesory.}
    \label{fig:enter-label}
\end{figure}

\begin{figure}[H]
    \centering
    \includegraphics[width=1\linewidth]{64_64_48.png}
    \caption{Rozmiar 64 na 64, 48 procesorów.}
    \label{fig:enter-label}
\end{figure}

\newpage

\subsection{Rozmiar 128 na 128}
\begin{table}[ht]
\centering
\begin{tabular}{|l|r|r|r|}
\hline
\rowcolor{orange!40} 
Liczba rdzeni & Średni czas wykonania algorytmu (s) \\\hline
1 (implementacja sekwencyjna) & 570.46 \\\hline
2 & 297.40 \\\hline
4 & 68.19 \\\hline
8 & 64.31 \\\hline
16 & 66.75 \\\hline
32 & 68.15 \\\hline
48 & 67.98 \\\hline
\end{tabular}
\caption{\label{tab:widgets}Średni czas obliczeń dla problemu o rozmiarze 128 na 128.}
\end{table}

\begin{figure}[H]
    \centering
    \includegraphics[width=1\linewidth]{128_128_4.png}
    \caption{Rozmiar 128 na 128, 4 procesory.}
    \label{fig:enter-label}
\end{figure}

\begin{figure}[H]
    \centering
    \includegraphics[width=1\linewidth]{128_128_8.png}
    \caption{Rozmiar 128 na 128, 8 procesorów.}
    \label{fig:enter-label}
\end{figure}

\begin{figure}[H]
    \centering
    \includegraphics[width=1\linewidth]{128_128_16.png}
    \caption{Rozmiar 128 na 128, 16 procesorów.}
    \label{fig:enter-label}
\end{figure}

\begin{figure}[H]
    \centering
    \includegraphics[width=1\linewidth]{128_128_32.png}
    \caption{Rozmiar 128 na 128, 32 procesory.}
    \label{fig:enter-label}
\end{figure}

\begin{figure}[H]
    \centering
    \includegraphics[width=1\linewidth]{128_128_48.png}
    \caption{Rozmiar 128 na 128, 48 procesorów.}
    \label{fig:enter-label}
\end{figure}

\section{Wnioski}
Ze względu na ograniczenia dostępu do większych ilości procesorów, udało się rozwiążać jedynie relatywnie małe problemy. Efektywności użycia CPU jest niska, ale widzimy, że wzrasta razem ze zwiększaniem rozmiaru problemu.

\end{document}